\section{Conclusions and Future Work}
\label{sec:conclusion}
This paper has presented a lightweight TypeScript-based 3D engine for browser-delivered interactive educational content, together with three optimization strategies addressing distinct dimensions of runtime performance. The engine wraps a low-level graphics library behind a high-level component model inspired by mature commercial engines, letting content teams build interactive 3D learning environments with familiar software-engineering patterns and without encountering the underlying graphics API.

The evaluation covered 87 configurations across three benchmark scenes, using three graphics adapters on two machines. Two results shape how the engine should develop. The scene with the highest draw-call count proved bounded by main-thread submission rather than by the GPU, with the graphics adapter on a fixed host leaving its render phase unchanged to within 0.14\%. And one strategy favoured the discrete adapters while one flag combination penalized only the integrated one, so the benefit of an optimization does not follow a hardware-tier ordering. Together these indicate that targeted improvements at the architecture and asset-pipeline levels can substantially reduce the startup latency of browser-based educational 3D applications without a native runtime exported to WebAssembly.

Future work follows from the measurements. The draw-call bound motivates exercising the engine's existing GPU instancing and static-batching facilities on representative educational content, which the benchmark scenes did not do. Extending the evaluation to mobile platforms remains outstanding, as does a WebGPU backend requiring a generalized adapter layer. A graphical scenario editor is planned, and opening the platform to students will enable the user study these controlled measurements cannot substitute for.